% =============================================================================
% 00_titelblatt.tex — Titelblatt der IDPA
% =============================================================================
% Dieses Blatt bildet die offizielle Titelseite gemäss Schulvorgabe.
% Ersetze alle Platzhalter (<<...>>) durch die tatsächlichen Angaben.
% =============================================================================

\begin{titlepage}
  \centering

  % --- Logo und Name der Schule ----------------------------------------------
  % Binde das offizielle Schullogo ein (Datei z. B. images/schullogo.png).
  % Falls kein Logo vorhanden ist, kann die Zeile \includegraphics auskommentiert
  % und durch den Schulnamen in Textform ersetzt werden.
  % Passe die Dateiendung an das tatsächliche Format an (.png, .pdf, .jpg)
  \includegraphics[width=5cm]{schullogo.png}\\[0.5cm]
  {\Large \textbf{<<Name der Schule>>}}\\[0.3cm]
  {\large <<Abteilung / Fachrichtung>>}

  \vspace{2cm}

  % --- Art der Arbeit --------------------------------------------------------
  {\large Interdisziplinäre Projektarbeit (IDPA)}\\[0.5cm]

  % --- Vollständiger Titel der Arbeit ----------------------------------------
  {\Huge \textbf{Geotracker}}\\[0.4cm]
  {\Large \textbf{<<Untertitel der Arbeit>>}}

  \vspace{2cm}

  % --- Autor:innen und Klasse ------------------------------------------------
  % Füge alle beteiligten Personen mit Klasse ein.
  \begin{tabular}{ll}
    \textbf{Autor:innen:} & <<Vorname Nachname>>, Klasse <<Klasse>> \\
                          & <<Vorname Nachname>>, Klasse <<Klasse>> \\
                          & <<Vorname Nachname>>, Klasse <<Klasse>> \\
  \end{tabular}

  \vspace{1.5cm}

  % --- Betreuende Lehrpersonen -----------------------------------------------
  \begin{tabular}{ll}
    \textbf{Betreuung:} & <<Titel Vorname Nachname>> \\
                        & <<Titel Vorname Nachname>> \\
  \end{tabular}

  \vspace{1.5cm}

  % --- Abgabedatum -----------------------------------------------------------
  \begin{tabular}{ll}
    \textbf{Abgabedatum:} & <<TT. Monat JJJJ>> \\
  \end{tabular}

  \vfill

  % --- Schuljahr -------------------------------------------------------------
  {\small Schuljahr <<JJJJ/JJJJ>>}

\end{titlepage}

% Seitenzähler zurücksetzen, damit das Titelblatt keine Seitenzahl trägt
\setcounter{page}{0}
